% TOP_PROBLEM: {"id":30007626,"problem_number":"LOCAL-30007626","title":"Costs of geopolitical fragmentation","category_id":21,"category":{"id":21,"name":"economics","display_name":"Economics","description":"Open economic theory statements, published research agendas, and proposed scoped specifications; question provenance and historical priority are qualified per record.","slug":"economics","order_index":21,"created_at":"2026-10-08T00:00:00Z"},"status":"research_question","external_url":null,"legacy_ids":[],"published":false,"statement":"In the explicitly specified setting: How large are the dynamic welfare and productivity costs of decoupling once firms adapt and trade is rerouted? The mathematical statement above fixes the target, assumptions and scope of this catalogue formulation.","economics":{"original_id":115,"field":"Trade, globalization and geoeconomics","kind":"empirical","formalization_basis":"adapted_research_question","source_status":"research_agenda","priority_status":"earliest_found","priority_note":"This is the earliest explicit statement verified in this audit; first priority for the full catalogue formulation is not established.","earliest_verified_reference":{"authors":["Bank of England"],"title":"Bank of England Agenda for Research, 2025–2028","year":2026,"url":"https://www.bankofengland.co.uk/research/bank-of-england-agenda-for-research","doi":null,"locator":"Interconnected World / Priority topics for2026, item4","evidence_summary":"Asks fragmentation effects on activity, inflation and stability."},"current_reference":{"authors":["Bank of England"],"title":"Bank of England Agenda for Research, 2025–2028","year":2026,"url":"https://www.bankofengland.co.uk/research/bank-of-england-agenda-for-research","doi":null,"locator":"Interconnected World / Priority topics for2026, item4","evidence_summary":"Asks fragmentation effects on activity, inflation and stability."},"formalization_note":"The source asks fragmentation effects broadly. Consumption-equivalent losses with endogenous firm adaptation form a proposed representative specification. Independent math audit: the consumption-equivalent factor preserves positive consumption, permits gains, and requires the comparison utility to be in the attainable range."},"statusQualification":"Published question or research agenda verified at the cited locator. The mathematical specification is adapted for this catalogue and includes additional assumptions; source publication does not establish first-ever priority or certify this exact model remains unsolved. The source asks fragmentation effects broadly. Consumption-equivalent losses with endogenous firm adaptation form a proposed representative specification. Independent math audit: the consumption-equivalent factor preserves positive consumption, permits gains, and requires the comparison utility to be in the attainable range.","statusReviewedAt":"2026-10-08"}
% FIRST_POSTED: 2026-10-08
% ENTRY_KIND: statement_only
% !TeX program = lualatex
\documentclass[11pt]{article}
\usepackage{fontspec}
\usepackage[a4paper,margin=25mm]{geometry}
\usepackage{amsmath,amssymb,mathtools}
\usepackage{xurl}
\usepackage[hidelinks]{hyperref}
\newcommand{\sref}[2]{\hyperlink{source:#1:#2}{[S#2]}}
\setlength{\emergencystretch}{3em}
\begin{document}
% problemId:30007626
\section*{Costs of geopolitical fragmentation}\label{30007626}
\subsection{Definitions and mathematical statement}
Fix a multi-country, multi-sector economy with trade and investment networks. A fragmentation scenario \(a\) specifies bilateral restrictions by partner, sector, asset type and implementation date. Firms may switch suppliers, enter, exit and reorganize production, subject to explicit adjustment technologies; knowledge diffusion and investment are included only through declared equations. Let \(C_{ct}(a)\) denote representative-country or distribution-weighted consumption in equilibrium. Define the discounted welfare loss using a consumption-equivalent factor \(\lambda_c\):
\[E\sum_{t\geq0}\beta^t u_c((1-\lambda_c)C_{ct}(0))=E\sum_{t\geq0}\beta^t u_c(C_{ct}(a)),\]
where \(\beta\in(0,1)\), strictly increasing utility \(u_c\) and the no-fragmentation path \(0\) are fixed, and sums converge. Require \(\lambda_c<1\) and assume the comparison utility lies in the range generated by proportional baseline consumption; negative values indicate a welfare gain. The task is to identify or bound \(\lambda_c\) and world productivity effects under plausible adaptations and trade rerouting. State restrictions identifying substitution elasticities, fixed reorganization costs and geopolitical selection separately. Report sensitivity to bloc definitions, connector countries and the probability of future policy reversal. Distinguish already observed partial fragmentation from hypothetical complete decoupling. Model-based scenario estimates answer conditional counterfactuals; they are not identified causal global losses merely because a calibrated model reproduces baseline trade shares.

\par\textbf{Formulation and scope.} The source asks fragmentation effects broadly. Consumption-equivalent losses with endogenous firm adaptation form a proposed representative specification. Independent math audit: the consumption-equivalent factor preserves positive consumption, permits gains, and requires the comparison utility to be in the attainable range.

\par\textbf{Open-question provenance.} An explicit question or research agenda was verified at the source locator. This does not by itself certify that every later version remains unresolved \sref{30007626}{1}.

\par\textbf{Historical priority.} This is the earliest explicit statement verified in this audit; first priority for the full catalogue formulation is not established.

\subsection{Short English statement}
In the explicitly specified setting: How large are the dynamic welfare and productivity costs of decoupling once firms adapt and trade is rerouted? The mathematical statement above fixes the target, assumptions and scope of this catalogue formulation.

\subsection{Sources}
\begin{itemize}\raggedright
\item[\textnormal{[S1]}]\hypertarget{source:30007626:1}{} Bank of England. 2026. Bank of England Agenda for Research, 2025–2028. Interconnected World / Priority topics for2026, item4.
\url{https://www.bankofengland.co.uk/research/bank-of-england-agenda-for-research}.
{\small\textit{Earliest verified question/agenda reference; historical priority qualified below. Asks fragmentation effects on activity, inflation and stability.}}
\end{itemize}
\end{document}
